% Figure: stress concentration factor for a filleted shaft step versus the
% ratio of fillet radius to small-shaft diameter, schematic curves for two
% D/d ratios. Shows how a larger fillet radius lowers Kt.
\begin{figure}[htbp]
\centering
\begin{tikzpicture}
\begin{axis}[
  width=12cm, height=7.5cm,
  xlabel={fillet radius ratio $r/d$},
  ylabel={stress concentration factor $K_t$},
  xmin=0, xmax=0.3, ymin=1, ymax=3.2,
  axis lines=left,
  tick label style={font=\scriptsize},
  label style={font=\small},
  legend pos=north east,
  legend style={font=\scriptsize},
]
% schematic Kt curves: decreasing with r/d, asymptotic. Use a fitted decaying form.
% Kt = 1 + A * (r/d)^(-b) bounded; use Kt = 1 + 1.9*exp(-9*r/d) for D/d=2 (bending)
\addplot[engblue, very thick, samples=100, domain=0.02:0.3]
  {1 + 1.9*exp(-9*x)};
\addlegendentry{$D/d=2$ (bending)}
\addplot[engred, very thick, samples=100, domain=0.02:0.3]
  {1 + 1.4*exp(-9*x)};
\addlegendentry{$D/d=1.5$ (bending)}
% the "keep below 2" design line
\addplot[enggray, dashed, samples=2, domain=0:0.3] {2};
\node[enggray, font=\scriptsize, anchor=south east] at (axis cs:0.3,2) {design target $K_t<2$};
\end{axis}
\end{tikzpicture}
\caption[Fillet stress concentration factor]{Schematic stress
concentration factor $K_t$ for a stepped shaft under bending, as a
function of the fillet radius normalised by the small-shaft diameter
$d$, for two step ratios $D/d$. A sharp shoulder ($r\to 0$) drives
$K_t$ toward the elastic singularity; a generous fillet brings it
toward unity. The standard design discipline keeps $K_t$ below about
$2$ at every geometric transition, which here means a fillet radius of
roughly $r/d>0.1$. The curves are schematic; design uses tabulated
charts such as Peterson's for the exact factor of a given geometry.}
\label{fig:vol03ch03:fillet-kt}
\end{figure}
